\subsection*{بخش اول}
\subsection*{الف}
با فرض $o_t=1$ و $f_t=1-i_t$ داریم:
\begin{flalign*}
h_t &= \tanh(c_t), \qquad
c_t=(1-i_t)\odot c_{t-1}+i_t\odot \tilde{c}_t,\\
\tilde{c}_t &= \tanh(W_c x_t+U_c h_{t-1}+b_c), \qquad
c_{t-1}=\tanh^{-1}(h_{t-1}) &&
\end{flalign*}
پس:
\begin{flalign*}
h_t=\tanh\!\Big((1-i_t)\odot \tanh^{-1}(h_{t-1})+i_t\odot \tanh(W_c x_t+U_c h_{t-1}+b_c)\Big). &&
\end{flalign*}

\subsection*{ب}
در \lr{GRU} داریم:
\begin{flalign*}
h_t=(1-z_t)\odot h_{t-1}+z_t\odot \tilde{h}_t. &&
\end{flalign*}
شباهت: $i_t$ نقش گیت به‌روزرسانی مانند $z_t$ دارد و بین حالت قبلی و کاندیدا وزن‌دهی می‌کند.  
تفاوت: در \lr{LSTM} ترکیب روی $c_{t-1}$ انجام شده و سپس $\tanh$ بیرونی اعمال می‌شود، و گیت بازنشانی $r_t$ وجود ندارد.

\subsection*{بخش دوم}
\subsection*{الف}
نقش $r_t$ در \lr{GRU} کنترل میزان استفاده از $h_{t-1}$ در ساخت کاندیدا است؛ $r_t\!\approx\!0$ یعنی قطع اثر گذشته. برای مشابه‌سازی در \lr{LSTM} می‌توان:
\begin{flalign*}
r_t&=\sigma(W_r x_t+U_r h_{t-1}+b_r),\\
\tilde{c}_t&=\tanh\!\big(W_c x_t+U_c(r_t\odot h_{t-1})+b_c\big). &&
\end{flalign*}

\subsection*{ب}
فرض $o_t=1$ یعنی خروجی مستقیماً برابر $\tanh(c_t)$ است و حذف کنترل خروجی؛
فرض $f_t=1-i_t$ یعنی یک گیت واحد نقش نگه‌داری/به‌روزرسانی را بازی می‌کند. با افزودن $r_t$، هر دو مدل دو گیت اصلی دارند: گیت به‌روزرسانی و گیت بازنشانی؛ تفاوت اصلی در وجود حالت سلولی $c_t$ و $\tanh$ بیرونی در \lr{LSTM} است.